% !TEX program = xelatex
\documentclass[UTF8,a4paper,11pt]{ctexart}

\usepackage[a4paper,margin=2.45cm,headheight=15pt]{geometry}
\usepackage{amsmath,amssymb}
\usepackage{booktabs,tabularx,array}
\usepackage{enumitem}
\usepackage{fancyhdr}
\usepackage[hidelinks]{hyperref}

\newcolumntype{L}[1]{>{\raggedright\arraybackslash}p{#1}}
\newcommand{\code}[1]{\texttt{\detokenize{#1}}}
\setlist[itemize]{topsep=3pt,itemsep=2pt,parsep=0pt}
\setlength{\parindent}{2em}
\setlength{\parskip}{0.35em}
\setlength{\emergencystretch}{2em}

\pagestyle{fancy}
\fancyhf{}
\fancyhead[L]{\small MultiGrid 汇报讲稿}
\fancyhead[R]{\small 张鑫}
\fancyfoot[C]{\thepage}
\renewcommand{\headrulewidth}{0.4pt}

\title{\bfseries High-Performance Implementation of\\
Multigrid Solver for Lattice QCD\\[0.4em]
\large 10 分钟汇报讲稿}
\author{张鑫\\中国科学院近代物理研究所}
\date{2026 年 10 月 11 日}

\begin{document}

\maketitle
\thispagestyle{empty}

\noindent\textbf{演讲总时长：}10 分钟。横版汇报稿共 17 页，建议平均每页
35 秒；公式和算法页约 35--40 秒，性能图页约 30--35 秒，封面和致谢页
各约 10 秒。以下文字为逐页口播稿，不是论文摘要；遇到时间不足时，
优先保留每页加粗的判断句和数字。

\begin{center}
\small
\begin{tabular}{@{}llll@{}}
\toprule
页码 & 主题 & 建议时长 & 必须保留的信息 \\
\midrule
1 & 标题与范围 & 15 s & 只汇报 MultiGrid，数据基线为 test27 \\
2--4 & Wilson/Clover/Dslash/Schur & 100 s & 传播子、Clover 块、奇偶消元 \\
5--8 & MultiGrid 与算法 & 140 s & 粗空间、Galerkin、Bi-CGStab、FGMRES \\
9 & CUDA--C++ 优化 & 40 s & 细层带宽、粗层同步、通信和显存 \\
10--14 & QUDA 对照与结论 & 175 s & 66/132、35/44、21/22、容量边界 \\
15--16 & 复现与先前成果 & 45 s & 最小闸门、508 报告的 64→1024 扩展 \\
17 & 致谢 & 10 s & 一个结论和一个行动项 \\
\bottomrule
\end{tabular}
\end{center}

\section*{逐页讲稿}

\subsection*{第 1 页｜标题与汇报范围（约 15 秒）}

各位老师、各位同行，大家好。我汇报的题目是面向格点 QCD 传播子求解的
高性能多重网格实现。这里重点讲 PyQCU 的 MultiGrid，尤其是 CUDA--C++
核心；性能数据统一采用 \code{git tag test27}，不把历史标签或实验性
加速比混入最终结论。

\emph{视觉提示：停在标题、项目和 test27 三个位置上，不展开项目历史。}

\subsection*{第 2 页｜Wilson 与 Clover 算子（约 35 秒）}

首先固定细网格算子。裸跃迁核 \(H\) 由四个正向和四个反向 neighbors 组成；
Wilson 算子是 \((m_0+4)I-\frac12H\)，等价地写成
\((m_0+4)(I-\kappa H)\)，其中
\(\kappa=1/(2m_0+8)\)。Clover 项是在相同最近邻结构上增加一个 onsite
的 \(12\times12\) spin--color 块 \(A_p=I+T_p\)。需要特别注意，C++ 的
原始 Dslash 入口返回 \(H\)，不是完整算子，也不包含 \(-\kappa\) 因子。

\emph{视觉提示：先指 Wilson 公式，再指 Clover onsite block，最后指
“raw H”警示栏。}

\subsection*{第 3 页｜原始 Dslash（约 30 秒）}

原始 Dslash 是矩阵自由的最近邻稀疏作用。每个格点读取八个邻居；对每个
方向执行 \(1\pm\gamma_\mu\) 投影后，再做规范平移和 \(SU(3)\) 矩阵乘法。
因此它天然是 bandwidth-bound。PyQCU 将四分量 spin projection 降到
两个独立分量，并利用 \(SU(3)\) 结构和时空维尾置布局减少访存。这里的
核心不是“多写几个 kernel”，而是把 Dslash 从条件分支和随机访存中解耦，
使其成为 GPU 擅长执行的连续数据流。

\emph{视觉提示：沿右侧五个流程节点快速移动，不逐行朗读入口表。}

\subsection*{第 4 页｜奇偶 Schur 消元（约 35 秒）}

为了减小外层 Krylov 维度，PyQCU 将场分成 even 和 odd 两个子空间。
以奇数子格为例，完整块算子包含 \(A_e,A_o\) 和两个 hopping block。
精确消去 even 分量后得到
\(S_o=A_o-\kappa^2H_{oe}A_e^{-1}H_{eo}\)。
求解 \(S_ox_o=b_o^{\mathrm{pc}}\) 后，再用
\(x_e=A_e^{-1}(b_e+\kappa H_{eo}x_o)\) 重建完整解。
这是精确代数变换，最终正确性仍由完整 Wilson/Clover 算子的真残差决定。

\emph{视觉提示：按“分块—消元—求解—重建”四步讲，避免陷入矩阵推导。}

\subsection*{第 5 页｜MultiGrid 的核心构造（约 35 秒）}

MultiGrid 的关键是把误差分解成高频和低频。平滑器压低高频误差，粗空间
表示低频误差。对每一层，先通过 \(A_lB_l\approx0\) 生成近零模，再在
aggregate 内做局部正交化，得到延拓 \(P_l\)，并令
\(R_l=P_l^\dagger\)。粗算子通过 Galerkin 投影
\(A_{l+1}=R_lA_lP_l\) 生成。局部正交意味着限制过程不需要全局 Gram
求逆，同时投影后的 coarse operator 仍保持有限支撑。

\emph{视觉提示：指公式 \(A_{l+1}=R_lA_lP_l\) 和右侧 V-cycle 流程。}

\subsection*{第 6 页｜预条件 Bi-CGStab（约 40 秒）}

这一页给出粗层求解器的完整算法。初始化残差后，每轮先计算
\(\rho_{i-1}\)，构造搜索方向 \(p^{(i)}\)；然后通过预条件器求
\(\hat p\)，应用原算子得到 \(v^{(i)}\)，再计算
\(\alpha_i\) 和中间残差 \(s\)。如果 \(s\) 已经满足收敛条件，就立即更新
\(x\) 并停止；否则继续求 \(\hat s\)、\(t=A\hat s\)，计算
\(\omega_i=t^\dagger s/(t^\dagger t)\)，完成最终的 \(r\) 和 \(x\) 更新。
粗层实现还保留热启动，但接近 breakdown 时回退冷启动。

\emph{视觉提示：只用手指沿表逐行扫过，不朗读所有公式；强调
\(M\hat p=p\) 和 \(\omega\) 两项。}

\subsection*{第 7 页｜PyQCU Strict 特化（约 35 秒）}

PyQCU Strict 将粗层分解放成 \(D_l=X_l+H_l\)，再做左预条件
\(\widehat D_l=X_l^{-1}D_l\)，最后投影得到
\(D_{l+1}=R_l\widehat D_lP_l\)。运行期不保存完整粗矩阵，而保存 onsite
\(X,X^{-1}\)、轴链接 \(Y\) 和预条件后的
\(\widehat Y=X^{-1}Y\)。细层是 compact target-parity 问题，粗层始终
保留 full geometry。这个顺序非常重要：必须先在 onsite 层做
\(X^{-1}\)，再做 \(P/R\)，否则会改变算子语义。

\emph{视觉提示：强调公式顺序和右侧 full coarse 数据流。}

\subsection*{第 8 页｜V-cycle 与 FGMRES（约 40 秒）}

一次 MG 预条件器在内部执行 pre-smoothing、restriction、递归粗层求解、
prolongation 和 post-smoothing。最粗层用 Bi-CGStab，中间层递归。
外层不直接用 CG，而是使用 flexible right-preconditioned GMRES：
每次 Arnoldi 列都允许使用不同的预条件器。右表给出 FGMRES 的完整循环，
包括 modified Gram--Schmidt、Givens 更新和 restart 后的 full true
residual 刷新。这里的迭代计数只记外层 FGMRES，不能把 smoother 和
coarse solver 的迭代相加。

\emph{视觉提示：先讲左侧递归，再讲右侧 FGMRES；最后指工作区公式。}

\subsection*{第 9 页｜CUDA--C++ 优化（约 40 秒）}

CUDA--C++ 的优化分为四层。细层 kernel 处理 spin projection、\(SU(3)\)
重用和合并访存；粗层通过 device-side scalars、融合标量更新和 CUDA Graph
减少 launch 与 host synchronization；setup 使用分块/着色 Galerkin、
显存预算、持久 hierarchy 和 SHA256 runtime cache；多 rank 通信使用
rank-local 粗 halo、非阻塞 vector halo 和远端 correction 融合。
complex64 默认启用 overlap，complex128 因浮点累加顺序敏感默认关闭。

\emph{视觉提示：按四列从上到下讲，不展开 API 和数组长度。}

\subsection*{第 10 页｜QUDA 对照口径（约 30 秒）}

性能对照严格采用 test27 的最终矩阵。请求矩阵是 72 个 combined 单元，
精确完成 66 个 combined、132 条 side record；其余 6 个双 P100 c64
大格请求项因显存或 570 秒 setup 时限被显式替代。每侧执行 1 次 cold、
2 次 warmup、5 次 steady，正式性能使用 trace-off。所有精确单元都
通过 config、input bundle 和 full-operator 真残差门。

\emph{视觉提示：四个 KPI 数字只需读出 72、66、132、6；右下角专门说明
“替代不是精确结果”。}

\subsection*{第 11 页｜总体分布（约 35 秒）}

定义 \(R=t_{\mathrm{QUDA}}/t_{\mathrm{PyQCU}}\)，\(R>1\) 表示 PyQCU
更快。66 个精确单元的整体中位比值是 1.025，说明总体只是接近持平。
但如果只看 MultiGrid，MG-2/3 有 35/44 个单元优于 QUDA；trace-off
MG 则有 21/22 个单元有利。相反，BiCGStab reference 是 0/22，
QUDA 全部更快。因此我们的结论必须限定在 MG 场景，不能外推到全部求解器。

\emph{视觉提示：先指整体中位，再指 MG 和 BiCGStab 两组，不逐点读数据。}

\subsection*{第 12 页｜设备与精度分组（约 35 秒）}

设备趋势非常清楚。V100 上 MG-3 的优势最大：c64 和 c128 的 trace-off
分组中位比分别是 4.74 和 3.22。双 P100 仍然总体有利，但幅度较小，
MG-2/3 都在 1.15 到 1.79 左右。图中同时显示极值：V100 c128 MG-3
最大单点达到 14.63，但也存在 0.54 的慢点。因此中位数表示趋势，
极值提醒我们结果对格点和残差轨迹敏感。

\emph{视觉提示：指色块从 V100 到 P100 的变化，不使用“全面领先”。}

\subsection*{第 13 页｜阶段、容量、残差与显存（约 40 秒）}

trace-on 阶段分解显示，PyQCU 的 coarse solve 与 other 合计约占 95\%；
这说明进一步优化细层算术的收益有限，重点应转到粗层同步、残差检查和
通信。容量方面，双 P100 c64 大格的三条路径分别约为 15.9、14.8 和
12.4 GB；CPU staging 虽降低显存，但 setup 仍超过 570 秒。因此 6 个
请求项被替换为 \(16^4\)，且不进入原大格速度比。正确性方面，PyQCU
和 QUDA 的最大真残差分别为 \(7.40\times10^{-7}\) 和
\(7.97\times10^{-7}\)，均在 \(5\times10^{-6}\) 门限以内。

\emph{视觉提示：左侧指阶段占比，右侧指容量替代表，底部指残差和显存 KPI。}

\subsection*{第 14 页｜结论、边界与下一步（约 35 秒）}

当前结论可以压缩成三句话。第一，MG 专项优势明确，尤其 V100 上显著；
第二，整体与 BiCGStab 仍暴露明显边界，trace-on 同步和 P100 大格容量
也限制推广；第三，下一步优先做粗层同步、持久 request、device-aware
通信、setup 显存和更多平台的端到端验证。我们不宣称任意格点、任意层数
或所有求解器全面领先。

\emph{视觉提示：三栏分别讲优势、不足、下一步，最后强调不可外推。}

\subsection*{第 15 页｜参考实现与复现（约 20 秒）}

参考实现包括 quantum-mg、QUDA 和 DD-alphaAMG。复现顺序是先构建，
再运行 Strict fast gate，随后运行双 rank MPI probe，最后才启动完整矩阵。
这样可以在较长 benchmark 之前先排除 ABI、生命周期、halo 和真残差问题。
当前平台优先级仍是 CUDA 性能优化，然后是 DCU/CANN 真机适配和
TileLang 内核化。

\emph{视觉提示：只读命令类别，不逐条念完整路径。}

\subsection*{第 16 页｜508 报告的大规模成果（约 25 秒）}

这一页补充先前成果。张鑫 508 报告在 4 卡 DCU 环境中记录了从 64 到
1024 进程规模的 Wilson 与 Clover 大格测试。报告中的累计加速比在
1024 进程时分别为 3.12 和 7.19，说明早期 DCU/CUDA 路径已经具备正确性
和大规模扩展趋势；右下图给出 Clover 与 Wilson 单次迭代耗时差随等效
体积增长的数据。需要明确，这是早期成果，当时 MG 仍属前期验证，不能
与 test27 的 Strict-MG 正式速度比混合。

\emph{视觉提示：先读三个 KPI，再指两条扩展曲线，最后指原始 TEST9 表。}

\subsection*{第 17 页｜致谢（约 10 秒）}

我的汇报到这里。核心结论是 MultiGrid 专项性能已经进入国际先进实现
同量级，并在 V100 等场景形成清晰优势；下一步要把粗层同步、通信和
setup 显存做成可复验的持续收益。谢谢大家。

\emph{视觉提示：停在总结句，随后进入提问。}

\section*{备用问答要点}

\begin{itemize}
  \item \textbf{为什么总中位只有 1.025？} 因为 66 个单元同时包含
        BiCGStab reference；MG 单独是 35/44，trace-off MG 为 21/22。
  \item \textbf{为什么 trace-on 不能作为正式性能？} trace-on 插入同步和
        记录，改变运行时开销；它只用于逐层归因，正式比值使用 trace-off。
  \item \textbf{为什么 c128 overlap 默认关闭？} 改变浮点累加顺序后，
        大格三层曾由 32 次外层迭代退化到 451 次，因此必须 fail-safe。
  \item \textbf{508 报告是否证明当前 Strict-MG 千卡扩展？} 不直接证明。
        它证明早期 DCU 传播子求解的大规模扩展能力；当前 Strict-MG 的
        公开强证据仍是 V100 单卡、双 P100 与 1/2/4 rank preflight。
  \item \textbf{最大的下一步是什么？} 降低 coarse solve 与同步成本，
        同时提升 setup 的显存效率，使未完成的大格请求进入正式精确矩阵。
\end{itemize}

\end{document}
